\chapter{The Proposed Method: \method}
\label{ch:method}

This chapter specifies \method{}, the Cross-Feature Fusion Mamba Network, and
its streaming variant \methodS{}. The still-image network fulfils the
supervised topic directly. It is a YOLO detector whose visible and thermal
features are fused by state-space blocks. The streaming variant adds the
contribution identified in Chapter~\ref{ch:novelty}: a causal state carried
from frame to frame. Every component has an ablation in
Chapter~\ref{ch:protocol}, and every claim is stated as a hypothesis in
Section~\ref{sec:hypotheses}.

\section{Design principles}
\label{sec:principles}

Five principles follow from the evidence of Chapters~\ref{ch:sources}
and~\ref{ch:sota}.

\begin{enumerate}[label=\textsc{p}\arabic*., leftmargin=2.6em]
  \item \emph{Keep what is known to work.} The base is YOLO26, with its
        small-target-aware assignment. A stride-4 head and stronger early
        stages are added, because those components produced the largest and
        cheapest gains in the source papers (Figure~\ref{fig:gains}).
  \item \emph{Spend state-space computation where sequences are long.} Within
        a frame, scans are applied at $P_3$ to $P_5$, where interleaving the
        two modalities doubles the sequence length. Across frames, the
        sequence is the stream itself and has no fixed length.
  \item \emph{Use time.} Memory carried across frames is the capability that
        none of the source papers, and none of the static fusion papers,
        provide.
  \item \emph{One control for three problems.} The discretisation step
        $\Delta$ is modulated by sensor reliability, by the true frame
        interval and by scene changes. In each case it decides how much the
        state should hold and how much it should write
        (Section~\ref{sec:delta}).
  \item \emph{Deployable by construction.} Every block has a pure-PyTorch
        path and a single-step form that updates the state once per frame,
        so that the streaming model can be exported and run on live
        cameras.
\end{enumerate}

\section{Architecture overview}
\label{sec:arch}

Figure~\ref{fig:arch} shows the full network for one frame $t$. Visible and
thermal frames enter two YOLO26 backbones. The backbones have separate
weights, because the two sensors have different statistics; sharing the
later stages is an ablation. Each produces features
$F^{v,l}_t$ and $F^{\theta,l}_t$ at levels $l\in\{2,3,4,5\}$. Here $v$ denotes
visible and $\theta$ thermal.

At $P_2$, fusion is a reliability-weighted sum, which is cheap enough for
160$\times$160 maps. At $P_3$ to $P_5$, the Cross-Modal Fusion Mamba (CMFM)
block of Section~\ref{sec:cmfm} produces fused features $Z^l_t$ and
reliability maps $\rho^l_t$. In the streaming variant, the Temporal State
Memory (TSM) of Section~\ref{sec:tsm} combines $Z^l_t$ with the state carried
from frame $t-1$. Before use, that state is warped by the estimated camera
motion $H_{t-1\to t}$. The fused, temporally enriched features enter a
YOLO26 path-aggregation neck extended to $P_2$. Heads at $P_2$, $P_3$ and
$P_4$ follow, with the $P_5$ head optional. Two side modules supply the
signals that control $\Delta$.
\begin{itemize}
  \item The \emph{stream monitor} reports the true interval $\tau_t$ since
        the last processed frame, sensor availability flags and a scene-cut
        test.
  \item The \emph{ego-motion estimator} computes $H_{t-1\to t}$.
\end{itemize}

\begin{figure}[p]
\centering
\begin{tikzpicture}[
    bv/.style={base, rgb, text width=1.55cm, minimum height=0.62cm, font=\scriptsize, inner ysep=2.5pt},
    bt/.style={base, ir,  text width=1.55cm, minimum height=0.62cm, font=\scriptsize, inner ysep=2.5pt},
    fz/.style={base, fuse, text width=1.75cm, minimum height=1.12cm, font=\scriptsize},
    tm/.style={base, tim, text width=1.75cm, minimum height=1.12cm, font=\scriptsize},
    nk/.style={base, n, text width=1.35cm, minimum height=1.12cm, font=\scriptsize},
    hdn/.style={base, fill=cFuse!9, text width=1.45cm, minimum height=1.12cm, font=\scriptsize},
    over/.style={preaction={draw=white, line width=3.2pt, -}},
  ]
  % columns
  \def\ca{0}\def\cb{2.25}\def\cc{4.75}\def\cd{7.35}\def\ce{9.85}\def\cf{12.05}
  % rows
  \def\yT{0.55}\def\yS{-1.0}\def\yb{-2.55}\def\yc{-4.1}\def\yd{-5.65}\def\ye{-7.2}\def\yo{-8.85}
  \def\bus{6.05}

  % ---------- column headings ----------
  \node[grp] at (\ca,1.75) {visible\\[-1pt]backbone};
  \node[grp] at (\cb,1.75) {thermal\\[-1pt]backbone};
  \node[grp] at (\cc,1.75) {cross-modal\\[-1pt]fusion};
  \node[grp] at (\cd,1.75) {temporal\\[-1pt]state memory};
  \node[grp] at (\ce,1.75) {neck};
  \node[grp] at (\cf,1.75) {heads};

  % ---------- inputs and side modules ----------
  \node[base, rgb, text width=1.75cm, minimum height=0.85cm, font=\scriptsize] (iv) at (\ca,\yT) {visible frame\\ $I^v_t$, $640{\times}512$};
  \node[base, ir,  text width=1.75cm, minimum height=0.85cm, font=\scriptsize] (it) at (\cb,\yT) {thermal frame\\ $I^\theta_t$, $640{\times}512$};
  \node[base, dat, text width=1.75cm, minimum height=0.85cm, font=\scriptsize] (mon) at (\cc,\yT) {stream monitor\\ $\tau_t$, $s^v_t$, $s^\theta_t$, cut};
  \node[base, n,   text width=1.75cm, minimum height=0.85cm, font=\scriptsize] (ego) at (\cd,\yT) {ego-motion\\ $H_{t-1\to t}$};
  \node[base, tim, fill=cTime!7, text width=2.95cm, minimum height=0.85cm, font=\scriptsize] (buf) at ({(\ce+\cf)/2},\yT) {state buffer $S^{3..5}_{t-1}$\\ written at frame $t-1$};

  % ---------- stems ----------
  \node[bv] (sv) at (\ca,\yS) {stem, $P_1$};
  \node[bt] (st) at (\cb,\yS) {stem, $P_1$};
  \node[base, tim, text width=1.75cm, minimum height=0.85cm, font=\scriptsize] (warp) at (\cd,\yS) {warp $\mathcal{W}$\\ $\hat S_{t-1}$ from $S_{t-1}$, $H$};
  \draw[arr] (iv) -- (sv);
  \draw[arr] (it) -- (st);
  \draw[arr] (ego) -- (warp);
  \draw[darr, draw=cTime] (buf.south) |- (warp.east);

  % ---------- backbone stages, staggered: visible high, thermal low ----------
  \foreach \l/\y/\res in {2/\yb/{160{\times}128}, 3/\yc/{80{\times}64}, 4/\yd/{40{\times}32}, 5/\ye/{20{\times}16}}{
    \node[bv] (v\l) at (\ca,{\y+0.28}) {$P_{\l}$\enspace$\res$};
    \node[bt] (t\l) at (\cb,{\y-0.28}) {$P_{\l}$\enspace$\res$};
  }
  \draw[arr] (sv) -- (v2);
  \draw[arr] (st) -- (t2);
  \foreach \i/\j in {2/3,3/4,4/5}{
    \draw[arr] (v\i) -- (v\j);
    \draw[arr] (t\i) -- (t\j);
  }

  % ---------- fusion column ----------
  \node[fz] (f2) at (\cc,\yb) {reliability-\\weighted sum};
  \foreach \l/\y in {3/\yc,4/\yd,5/\ye}
    \node[fz] (f\l) at (\cc,\y) {CMFM\\ $Z^{\l}_t$,\; $\rho^{\l}_t$};

  \foreach \l/\y in {2/\yb,3/\yc,4/\yd,5/\ye}{
    \draw[arr, over] (v\l.east) -- (v\l.east -| f\l.west);
    \draw[arr] (t\l.east) -- (t\l.east -| f\l.west);
  }
  \draw[darr, draw=cData] (mon) -- node[lbl, right=1pt, pos=0.45] {$s^m_t$} (f2);

  % ---------- temporal state memory ----------
  \foreach \l/\y in {3/\yc,4/\yd,5/\ye}
    \node[tm] (m\l) at (\cd,\y) {TSM\\ $\hat S^{\l}_{t-1}\to S^{\l}_t$};
  % state bus from the warp node to every TSM block
  \draw[darr, draw=cTime, line width=0.6pt] (warp.west) -| (\bus,\ye);
  \foreach \l/\y in {3/\yc,4/\yd,5/\ye}
    \draw[arr, draw=cTime] (\bus,{\y+0.3}) -- (\bus,{\y+0.3} -| m\l.west);
  \foreach \l/\y in {3,4,5}
    \draw[arr, over] (f\l.east) ++(0,-0.18) -- ($(m\l.west)+(0,-0.18)$);
  \draw[darr, draw=cData] (mon.east) -- node[lbl, above=0.5pt, pos=0.45] {$\tau_t$} (ego.west);

  % ---------- neck ----------
  \foreach \l/\y in {2/\yb,3/\yc,4/\yd,5/\ye}
    \node[nk] (n\l) at (\ce,\y) {PAN\\ level $P_{\l}$};
  \draw[arr, over] (f2.east) -- node[lbl, above=0.5pt, pos=0.58, fill=white] {no state at $P_2$} (n2.west);
  \foreach \l in {3,4,5}
    \draw[arr] (m\l) -- (n\l);
  \foreach \i/\j in {2/3,3/4,4/5}{
    \draw[arr] ($(n\j.north)+(-0.3,0)$) -- ($(n\i.south)+(-0.3,0)$);
    \draw[arr] ($(n\i.south)+(0.3,0)$) -- ($(n\j.north)+(0.3,0)$);
  }
  \node[lbl, anchor=south] at (\ce,{\yb+0.6}) {top-down $\uparrow$\enspace bottom-up $\downarrow$};

  % ---------- heads ----------
  \foreach \l/\y/\px in {2/\yb/4, 3/\yc/8, 4/\yd/16}
    \node[hdn, fill=cFuse!16] (h\l) at (\cf,\y) {head $P_{\l}$\\ stride \px};
  \node[hdn, fill=cFuse!6, text=muted] (h5) at (\cf,\ye) {head $P_5$\\ optional};
  \foreach \l in {2,3,4}
    \draw[arr] (n\l) -- (h\l);
  \draw[darr] (n5) -- (h5);

  % ---------- outputs ----------
  \node[base, fuse, text width=2.0cm, minimum height=0.95cm, font=\scriptsize] (out) at (\cf,\yo)
       {boxes, classes and\\ scores at frame $t$};
  \node[base, tim, fill=cTime!7, text width=1.95cm, minimum height=0.95cm, font=\scriptsize] (sto) at (\cd,\yo)
       {$S^{3..5}_t$ stored,\\ used at $t+1$};
  \draw[arr] (h5.south) ++(0,0) -- node[lbl, right=1pt] {all heads, NMS-free} (out);
  \draw[darr, draw=cTime] (m5) -- (sto);
  \node[lbl, text width=4.0cm, anchor=west] at (-0.95,\yo) {image mode: empty state\\ stream mode: state persists};
\end{tikzpicture}

\vspace{4mm}
\caption[Architecture of \method{} and \methodS{}]{Architecture of
\method{} and \methodS{} for one frame of a $640\times512$ visible--thermal
stream. Two YOLO26 backbones produce features at strides 4 to 32. A
reliability-weighted sum fuses $P_2$. CMFM blocks fuse $P_3$ to $P_5$ and
emit reliability maps $\rho$. In stream mode, the Temporal State Memory
updates a per-location state carried from frame $t-1$ after warping by the
estimated camera motion. Purple dashed paths carry state. Ochre dashed
paths carry the monitor's signals: sensor flags $s^m_t$, interval $\tau_t$
and scene cuts. The neck and heads follow YOLO26 with a stride-4 extension.
In image mode the state is empty and the network reduces to the still-image
\method{}.}
\label{fig:arch}
\end{figure}


\section{The small-object pathway}
\label{sec:smallpath}

Four measures protect small objects before any state-space block is
involved.
\begin{itemize}
  \item \emph{Stride-4 head.} The $P_2$ head is added and the neck carries
        $P_2$ through both its paths. In UAVDet and HPS-DETR this was the
        largest single source of small-object gain.
  \item \emph{Early capacity.} Channels and blocks are moved towards the
        first two stages, following UAVDet's ablation. The backbone
        configuration is a hyperparameter with three settings in the
        ablation.
  \item \emph{Native resolution.} RGBT-Tiny frames are $640\times512$ and
        are processed at native resolution with letterbox padding, so no
        downscaling destroys tiny objects. For larger inputs, sliced
        inference with SAHI~\citep{akyon2022sahi} is an inference option,
        not a training change.
  \item \emph{Assignment.} YOLO26's small-target-aware label assignment is
        retained~\citep{yolo26_2026}. The scale-adaptive SAFit measure
        proposed with RGBT-Tiny is used in evaluation~\citep{ying2025rgbttiny}.
\end{itemize}

\section{Cross-Modal Fusion Mamba with reliability-gated steps}
\label{sec:cmfm}

The CMFM block fuses $F^v, F^\theta \in \R^{C\times H\times W}$ at one pyramid
level. Superscript $l$ and subscript $t$ are dropped for clarity.
Figure~\ref{fig:cmfm} shows its structure. It has five steps.

\paragraph{Reliability estimation}
A small convolutional head $g$ estimates, at every location, how far each
modality can be trusted:
\begin{equation}
  \bigl(r^v,\, r^\theta\bigr) = \sigma\Bigl(g\bigl([\,F^v,\; F^\theta,\; |F^v-F^\theta|\,]\bigr)\Bigr)
  \odot \bigl(s^v,\, s^\theta\bigr),
  \qquad r^v, r^\theta \in (0,1)^{H\times W}.
  \label{eq:rel}
\end{equation}
Here $g$ is a depthwise $3\times3$ convolution followed by a $1\times1$
convolution with two outputs. The flags $s^v, s^\theta\in\{0,1\}$ come from
the stream monitor. They are zero when a sensor delivers no frame, so that a
known outage needs no learning. The absolute difference is included because
disagreement between modalities is the most direct evidence that one of
them is degraded.

\paragraph{Offset-guided alignment}
A second head predicts a displacement field $o \in \R^{2\times H\times W}$ from
$[F^v, F^\theta]$. The thermal features are resampled onto the visible grid,
$\tilde F^\theta = \mathcal{S}(F^\theta, o)$, with bilinear sampling. This follows
the offset guidance of COMO~\citep{como2024}. The field is initialised to
zero, so the block starts as an aligned fusion. Its effect is measured under
synthetic misalignment (Section~\ref{sec:robustness}).

\paragraph{Token interleaving}
Both streams are projected to $C'$ channels by $1\times1$ convolutions
$\phi_v$ and $\phi_\theta$. They are flattened along the scan order and
interleaved, so that the two modalities at each location are adjacent in
the sequence:
\begin{equation}
  X = \bigl(x^v_1,\; x^\theta_1,\; x^v_2,\; x^\theta_2,\; \dots,\; x^v_{HW},\; x^\theta_{HW}\bigr)
  \in \R^{2HW \times C'} .
  \label{eq:interleave}
\end{equation}
The hidden state therefore passes from one modality to the other at every
location. Fusion happens inside the state, not after it.

\paragraph{Reliability-gated selective scan}
The interleaved sequence is processed by a selective scan in the SS2D
pattern~\citep{liu2024vmamba}, with one change. The step of each token is
scaled by the reliability of its modality at its location:
\begin{equation}
  \Delta_k = \softplus\bigl(w_\Delta^{\top} x_k + b_\Delta\bigr)\cdot r^{m(k)}_{p(k)},
  \label{eq:gated-delta}
\end{equation}
where $m(k)\in\{v,\theta\}$ and $p(k)$ are the modality and location of token
$k$. By the argument of Section~\ref{sec:delta}, a token judged unreliable
receives $\Delta_k\to0$. It neither writes into the state nor erases it.

This differs in kind from the output gating of
Fusion-Mamba and CFMW~\citep{dong2025fusionmamba,li2025cfmw}. Output gating
suppresses a modality's contribution after it has already updated the state.
Gating the step prevents a corrupted token from contaminating the state in
the first place. The difference is small on still images. It becomes
decisive once the same state persists across frames.

\paragraph{Merge}
The scan output is de-interleaved into $Y^v$ and $Y^\theta$. A local branch,
a depthwise $3\times3$ convolution on $[\phi_v(F^v), \phi_\theta(\tilde F^\theta)]$,
runs alongside. This follows MambaVision's finding that a non-SSM branch
beside the scan helps~\citep{hatamizadeh2025mambavision}. The fused output
is
\begin{equation}
  Z = \hat r^v \odot F^v + \hat r^\theta \odot \tilde F^\theta
      + W_o\bigl[\,Y^v;\; Y^\theta;\; \text{local}\,\bigr],
  \qquad
  \hat r^m = \frac{r^m}{r^v + r^\theta + \epsilon},
  \label{eq:merge}
\end{equation}
which reduces to a reliability-weighted average when the scan contributes
nothing. The block also outputs the combined reliability
$\rho = \max(r^v, r^\theta)$, which the TSM uses.

\paragraph{Cost}
With $d$ scan directions and state size $N$, the scan costs
$\mathcal{O}(d\cdot 2HW\cdot C' N)$. At $P_3$ for a $640\times512$ input,
$2HW = 10\,240$ tokens, which is well inside the regime where a linear scan
is cheaper than attention. The number of directions ($d\in\{2,4\}$) is an
ablation.

\begin{figure}[htbp]
\centering
\begin{tikzpicture}[
    st/.style={base, text width=1.42cm, minimum height=1.0cm, font=\scriptsize},
    sm/.style={base, text width=1.1cm, minimum height=0.58cm, font=\scriptsize, inner ysep=2pt},
    over/.style={preaction={draw=white, line width=3.2pt, -}},
    tok/.style={circle, draw=none, inner sep=0pt, minimum size=1.55mm},
  ]
  \def\xa{0}\def\xb{1.9}\def\xc{3.75}\def\xd{5.55}\def\xe{7.4}\def\xf{9.3}\def\xg{11.2}\def\xh{12.8}
  \def\yv{1.3}\def\yt{-1.3}\def\yu{2.55}\def\yl{-2.45}

  % inputs
  \node[st, rgb] (fv) at (\xa,\yv) {$F^v$\\ $C{\times}H{\times}W$};
  \node[st, ir]  (ft) at (\xa,\yt) {$F^\theta$\\ $C{\times}H{\times}W$};

  % reliability and offset heads
  \node[st, n, text width=1.5cm, minimum height=1.35cm] (hd) at (\xb,0)
       {heads $g$, $h$\\ on $[F^v,F^\theta,$\\ $|F^v{-}F^\theta|]$};
  \draw[arr] (fv.south) to[out=-80,in=170] (hd.west |- 0,0.3);
  \draw[arr] (ft.north) to[out=80,in=-170] (hd.west |- 0,-0.3);

  % resample thermal onto the visible grid
  \node[st, ir] (rs) at (\xc,\yt) {resample\\ $\tilde F^\theta{=}\mathcal{S}(F^\theta\!,o)$};
  \draw[arr] (ft) -- (rs);
  \draw[arr] (hd.south east) -- node[lbl, right=0.5pt, pos=0.4] {$o$} (rs.north);

  % projections
  \node[sm, rgb] (pv) at (\xd,0.62) {$\phi_v$, $1{\times}1$};
  \node[sm, ir]  (pt) at (\xd,-0.62) {$\phi_\theta$, $1{\times}1$};

  % residual lanes (visible on top, thermal at the bottom)
  \draw[line width=0.55pt, draw=cNeutral] (fv.east) -- (\xd,\yv);
  \draw[arr] (\xd,\yv) -- (pv.north);
  \draw[line width=0.55pt, draw=cNeutral] (rs.east) -- (\xd,\yt);
  \draw[arr] (\xd,\yt) -- (pt.south);

  % interleave tokens
  \node[st, fuse, minimum height=1.2cm] (il) at (\xe,0) {interleave\\[13pt] \ };
  \foreach \i in {0,...,7}{
    \pgfmathsetmacro{\xx}{\xe-0.53+\i*0.152}
    \ifodd\i \node[tok, fill=cIR] at (\xx,-0.2) {}; \else \node[tok, fill=cRGB] at (\xx,-0.2) {}; \fi
  }
  \node[font=\tiny, text=muted] at (\xe,-0.43) {$x^v_1\,x^\theta_1\,x^v_2\,x^\theta_2\cdots$};
  \draw[arr] (pv.east) -- (il.west |- pv.east);
  \draw[arr] (pt.east) -- (il.west |- pt.east);

  % gated selective scan
  \node[st, tim, text width=1.55cm, minimum height=1.45cm] (sc) at (\xf,0)
       {selective scan\\ SS2D, $d$ dirs\\[1pt] $\Delta_k\cdot r^{m(k)}_{p(k)}$};
  \draw[arr] (il) -- (sc);

  % local branch
  \node[st, n] (lc) at (\xf,\yl) {local branch\\ DW $3{\times}3$};
  \draw[arr] (il.south) |- node[lbl, above=0.5pt, pos=0.75] {concat} (lc.west);

  % merge
  \node[st, fuse, text width=1.5cm, minimum height=1.45cm] (mg) at (\xg,0)
       {de-interleave\\ $[Y^v;Y^\theta;$ local$]$\\ $W_o$, residual};
  \draw[arr] (sc) -- (mg);
  \draw[arr] (lc.east) -| ($(mg.south)+(0.5,0)$);

  % residual paths around the scan
  \draw[arr] (\xd,\yv) -- (\xg,\yv) -- (mg.north);
  \draw[arr] (\xd,\yt) -- (\xg,\yt) -- (mg.south);
  \node[lbl, anchor=south, text=cRGB] at (\xe,\yv) {residual $\hat r^v\odot F^v$};
  \node[lbl, anchor=north, text=cIR] at (\xe,\yt) {residual $\hat r^\theta\odot\tilde F^\theta$};

  % reliability gate path along the top
  \draw[line width=0.6pt, draw=cData, dash pattern=on 2.2pt off 1.6pt] (hd.north) -- (\xb,\yu) -- (\xf,\yu);
  \draw[arr, over, draw=cData, dash pattern=on 2.2pt off 1.6pt] (\xf,\yu) -- (sc.north);
  \node[base, tim, fill=cTime!7, text width=1.95cm, minimum height=0.75cm, font=\scriptsize] (rho) at (\xg,\yu) {$\rho=\max(r^v\!,r^\theta)$};
  \draw[darr, draw=cData] (\xf,\yu) -- (rho.west);
  \node[lbl, anchor=south] at ({(\xb+\xf)/2},\yu) {$r^v, r^\theta \in (0,1)^{H\times W}$, multiplied by sensor flags $s^m$, gate the step of every token};

  % outputs
  \node[st, fuse, text width=1.05cm] (zo) at (\xh,0) {$Z$\\ to TSM\\ or neck};
  \draw[arr] (mg) -- (zo);
  \node[st, tim, fill=cTime!7, text width=1.05cm, minimum height=0.75cm] (ro) at (\xh,\yu) {to TSM};
  \draw[arr] (rho) -- (ro);
\end{tikzpicture}

\vspace{2mm}
\caption[The Cross-Modal Fusion Mamba block]{The Cross-Modal Fusion Mamba
(CMFM) block at one pyramid level. Small heads estimate per-location
reliability for each modality and a displacement field that resamples the
thermal features onto the visible grid. The projected streams are
interleaved so that the hidden state alternates between modalities at every
location. The selective scan multiplies each token's step $\Delta_k$ by the
reliability of its modality and location (ochre path). An unreliable token
therefore neither writes into the state nor erases it. A local
convolutional branch and reliability-weighted residuals complete the block.
The combined reliability $\rho$ is passed to the Temporal State Memory.}
\label{fig:cmfm}
\end{figure}


\section{Temporal State Memory}
\label{sec:tsm}

The TSM gives each location of each fused level a small state that persists
across frames. For level $l$ with resolution $H\times W$, the state is
$S_t \in \R^{H\times W\times C_s\times N}$, with $C_s=C/2$ and $N=8$ by
default.

\paragraph{Input}
The fused features are projected, $u_t = W_{\text{in}} Z_t \in \R^{H\times W\times C_s}$.
The input-dependent projections $B_t = W_B u_t$ and $C_t = W_C u_t$ are formed
as in Equation~\eqref{eq:selective}.

\paragraph{Motion compensation}
The previous state is warped into the current frame,
\begin{equation}
  \hat S_{t-1} = \mathcal{W}\bigl(S_{t-1};\, H_{t-1\to t}\bigr),
  \label{eq:warp}
\end{equation}
by bilinear sampling at the level's resolution. Locations that enter the
view have no history and receive a zero state.

\paragraph{Gated step}
The temporal step combines the learned selection with the three control
signals:
\begin{equation}
  \Delta_t(p) = \softplus\bigl(w^{\top} u_t(p) + b\bigr)\cdot
  \frac{\tau_t}{\tau_0}\cdot \rho_t(p).
  \label{eq:tsm-delta}
\end{equation}
Here $\tau_t$ is the measured interval since the previous processed frame
and $\tau_0$ the nominal interval seen in training. The factor $\tau_t/\tau_0$
keeps the discrete update faithful to the same continuous dynamics when
frames are dropped or arrive irregularly. This is the property Zubi\'c and
colleagues exploited for event cameras~\citep{zubic2024ssmevent}. The factor
$\rho_t$ holds the state where neither sensor can be trusted, so the memory
of an object survives a gap in the feed.

\paragraph{Update and readout}
For each location $p$,
\begin{align}
  S_t(p) &= \exp\bigl(\Delta_t(p)\,A\bigr)\odot \hat S_{t-1}(p)
            + \Delta_t(p)\; u_t(p)\, B_t(p)^{\top}, \label{eq:tsm-update}\\
  y_t(p) &= S_t(p)\, C_t(p) + D\odot u_t(p), \label{eq:tsm-read}\\
  \tilde Z_t &= Z_t + W_{\text{out}}\bigl(\mathrm{DWConv}_{3\times3}(y_t)\bigr).
  \label{eq:tsm-out}
\end{align}
The depthwise convolution in Equation~\eqref{eq:tsm-out} spreads temporal
evidence to neighbouring locations. That compensates for small errors in
the warp. Because the TSM is residual and the state starts at zero,
image mode is the special case of a single step from an empty state.
\method{} and \methodS{} therefore share one set of weights.

\paragraph{Reset}
The state is cleared in three cases:
\begin{itemize}
  \item the stream monitor detects a scene cut;
  \item motion estimation fails its inlier test;
  \item the gap $\tau_t$ exceeds a maximum $\tau_{\max}$.
\end{itemize}
A soft reset, which sets $\Delta$ large, is an alternative evaluated in the
ablation.

\paragraph{Cost}
Memory and computation per frame are $\mathcal{O}(HW\,C_s N)$ and do not depend
on how many frames have been seen. At $P_3$ for a $640\times512$ input with
$C_s=64$ and $N=8$, the state holds $80\times64\times64\times8 \approx 2.6$
million values, about 10\,MB in 32-bit floats. All three levels together
need about 14\,MB. Figure~\ref{fig:tsm} shows the memory unrolled over three
frames, including a thermal dropout and a frame gap.

\begin{figure}[htbp]
\centering
\begin{tikzpicture}[
    cl/.style={base, text width=2.4cm, minimum height=1.05cm, font=\scriptsize},
    wp/.style={op, fill=cTime!16, minimum size=6.2mm, font=\scriptsize},
  ]
  \def\dx{3.42}
  \def\yo{2.25}\def\ym{0.6}\def\yz{-1.05}\def\ys{-2.65}

  % frame headings
  \foreach \i/\lab in {0/{frame $t-1$}, 1/{frame $t$}, 2/{frame $t+1$}, 3/{frame $t+2$}}
    \node[grp] at (\i*\dx,3.35) {\lab};

  % --- frame t-1: normal ---
  \node[cl, dat]  (s0) at (0*\dx,\ys) {both sensors on\\ interval $\tau_0$};
  \node[cl, fuse] (z0) at (0*\dx,\yz) {$Z_{t-1}$\\ reliability $\rho\approx1$};
  \node[cl, tim]  (m0) at (0*\dx,\ym) {normal step\\ write and keep};
  \node[cl, fill=cFuse!8] (o0) at (0*\dx,\yo) {faint object\\ detected};

  % --- frame t: thermal dropout ---
  \node[cl, alert] (s1) at (1*\dx,\ys) {thermal dropout\\ flag $s^\theta_t=0$};
  \node[cl, fuse]  (z1) at (1*\dx,\yz) {$Z_t$ from visible\\ only; $\rho$ low at night};
  \node[cl, tim]   (m1) at (1*\dx,\ym) {$\Delta\to0$ where $\rho$\\ is low: state held};
  \node[cl, fill=cFuse!8] (o1) at (1*\dx,\yo) {object kept\\ from memory};

  % --- frame t+1: frame gap ---
  \node[cl, dat]  (s2) at (2*\dx,\ys) {two frames lost\\ $\tau_{t+1}=3\tau_0$};
  \node[cl, fuse] (z2) at (2*\dx,\yz) {$Z_{t+1}$\\ reliability $\rho\approx1$};
  \node[cl, tim]  (m2) at (2*\dx,\ym) {step scaled by\\ $\tau/\tau_0=3$};
  \node[cl, fill=cFuse!8] (o2) at (2*\dx,\yo) {position\\ re-anchored};

  % --- frame t+2: scene cut ---
  \node[cl, alert] (s3) at (3*\dx,\ys) {scene cut found\\ by the monitor};
  \node[cl, fuse]  (z3) at (3*\dx,\yz) {$Z_{t+2}$};
  \node[cl, tim, fill=cTime!6] (m3) at (3*\dx,\ym) {reset $S=0$\\ fresh start};
  \node[cl, fill=cFuse!8] (o3) at (3*\dx,\yo) {behaves as\\ image mode};

  % vertical flow inside each frame
  \foreach \i in {0,1,2,3}{
    \draw[arr] (s\i) -- (z\i);
    \draw[arr] (z\i) -- (m\i);
    \draw[arr] (m\i) -- (o\i);
  }

  % state carried between frames through the warp
  \node[wp] (w01) at (0.5*\dx,\ym) {$\mathcal{W}$};
  \node[wp] (w12) at (1.5*\dx,\ym) {$\mathcal{W}$};
  \draw[tarr] (m0) -- (w01); \draw[tarr] (w01) -- (m1);
  \draw[tarr] (m1) -- (w12); \draw[tarr] (w12) -- (m2);
  \node[lbl, anchor=south] at (w01.north) {$H_{t-1\to t}$};
  \node[lbl, anchor=south] at (w12.north) {$H_{t\to t+1}$};

  % reset between t+1 and t+2
  \node[op, fill=cAlert!14, minimum size=6.2mm] (cut) at (2.5*\dx,\ym) {$\times$};
  \draw[darr, draw=cTime] (m2) -- (cut);
  \node[lbl, anchor=south] at (cut.north) {state discarded};

  % incoming and outgoing state
  \draw[darr, draw=cTime] ($(m0.west)+(-0.85,0)$) -- node[lbl, above=0.5pt] {$S_{t-2}$} (m0.west);
  \draw[darr, draw=cTime] (m3.east) -- ++(0.75,0) node[lbl, above=0.5pt, pos=0.5] {$S_{t+2}$};

  % equation strip
  \node[font=\footnotesize, text=ink] at (1.5*\dx,-3.75)
    {$S_t = e^{\Delta_t A}\odot\mathcal{W}(S_{t-1};H)+\Delta_t\,u_t B_t^{\top},
      \qquad \Delta_t=\softplus(\cdot)\cdot\dfrac{\tau_t}{\tau_0}\cdot\rho_t$};
\end{tikzpicture}

\vspace{2mm}
\caption[The Temporal State Memory unrolled over four frames]{The Temporal
State Memory unrolled over four frames that exercise its three controls. At
$t-1$ both sensors work and the state updates normally. At $t$ the thermal
sensor drops out at night. The reliability $\rho$ falls where only thermal
could see, the step $\Delta$ falls with it, and the state holds the object
until the sensor returns. At $t+1$ two frames have been lost, so the step is
scaled by the true interval. At $t+2$ the monitor detects a scene cut and
the state is reset. Purple arrows carry the state through the motion warp
$\mathcal{W}$.}
\label{fig:tsm}
\end{figure}


\section{Ego-motion compensation and scene changes}
\label{sec:egomotion}

Drone and handheld footage moves. Without compensation, the state at a
location would describe whatever was there before the camera moved.
\begin{itemize}
  \item The default estimator is classical. ORB features are matched between
        consecutive frames at half resolution, and a homography is fitted
        with RANSAC. The thermal frame is used at night and the visible frame
        by day, chosen by the stream monitor from the global reliability. The
        result is scaled to each pyramid level.
  \item For static cameras the estimator is bypassed and $H$ is the
        identity.
  \item A learned alternative regresses the eight homography parameters
        from the $P_5$ features of frames $t-1$ and $t$. It is cheaper on an
        embedded device and is compared in the ablation.
\end{itemize}
A homography models only a planar or distant scene. Parallax and
independently moving objects are left to the learned update, the
depthwise spreading in Equation~\eqref{eq:tsm-out} and the tracker.

\section{Training}
\label{sec:training}

Training proceeds in two stages (Figure~\ref{fig:training}).

\paragraph{Stage A, still images}
\begin{itemize}
  \item Both backbones start from COCO-pretrained YOLO26 weights. The first
        convolution of the thermal backbone is initialised by averaging the
        pretrained filters over the colour channels.
  \item The CMFM blocks, neck and heads are trained on single frames from
        the visible--thermal datasets of Chapter~\ref{ch:data}.
  \item The result is the still-image \method{}, which is also the Phase~1
        deliverable.
\end{itemize}

\paragraph{Stage B, clips}
\begin{itemize}
  \item Training continues on clips of $L$ consecutive frames, with $L=8$
        and later $L=16$.
  \item Each clip uses a random frame stride $s\in\{1,2,3\}$, so that the
        model sees varied intervals $\tau_t$.
  \item The first two frames of each clip are a burn-in without loss.
        Gradients flow through the clip and are truncated at its boundary
        (truncated back-propagation through time~\citep{williams1990tbptt}).
  \item With probability 0.5, the final state of one clip initialises the
        next clip of the same sequence. The model therefore experiences
        states that are much older than $L$ frames.
\end{itemize}

\paragraph{Augmentation and corruption}
Geometric augmentation is identical for both modalities and for every frame
of a clip. Mosaic is disabled in Stage B because it breaks temporal
continuity. Photometric jitter is applied to the visible stream only.
Thermal frames receive gain and offset drift and sensor noise. A corruption
simulator then applies five perturbations per clip.
\begin{itemize}
  \item Bursty dropout of each modality is drawn from a two-state Markov
        chain, with probability 0.02 of failing and 0.2 of recovering per
        frame. Outages therefore last five frames on average.
  \item Rare total blackouts remove both sensors at once.
  \item Thermal misalignment is a random affine offset, up to 16 pixels of
        translation, $\pm3^\circ$ of rotation and $\pm5\%$ of scale. It
        drifts slowly within the clip.
  \item Desynchronisation delays the thermal stream by zero to two frames.
  \item Random frame drops make the intervals irregular.
\end{itemize}
Because the simulator knows what it corrupted, it also provides targets for
the reliability heads.

\paragraph{Objective}
The loss averages the YOLO26 detection loss over the supervised frames of a
clip, $\mathcal{T}$, and adds a weak reliability term:
\begin{equation}
  \mathcal{L} = \frac{1}{|\mathcal{T}|}\sum_{t\in\mathcal{T}} \mathcal{L}_{\text{det}}(t)
  \;+\; \lambda_{\text{rel}}\,\frac{1}{|\mathcal{T}|}\sum_{t\in\mathcal{T}}\;\sum_{m\in\{v,\theta\}}
  \mathrm{BCE}\bigl(\bar r^m_t,\; q^m_t\bigr),
  \qquad \lambda_{\text{rel}} = 0.1 .
  \label{eq:loss}
\end{equation}
Here $\bar r^m_t$ is the spatial mean of a modality's reliability map and
$q^m_t\in\{0,1\}$ is the simulator's target. Optimisation uses the YOLO26
defaults with mixed precision. The selective scan and the TSM state are kept
in 32-bit precision, because T4 GPUs lack bfloat16 and the scan is sensitive
to half-precision rounding. An exponential moving average of the weights is
kept for evaluation.

\begin{figure}[htbp]
\centering
\begin{tikzpicture}[
    tn/.style={base, text width=1.42cm, minimum height=1.5cm, font=\scriptsize},
    over/.style={preaction={draw=white, line width=3.2pt, -}},
  ]
  \def\dx{2.04}
  \def\ya{0}\def\yb{-3.35}\def\yc{-7.15}

  % ---------- stage A ----------
  \node[grp, anchor=west] at (-0.8,{\ya+1.2}) {Stage A, still images};
  \node[tn, dat]  (a0) at (0*\dx,\ya) {still pairs:\\ LLVIP, M3FD,\\ FLIR and\\ RGBT-Tiny};
  \node[tn, n]    (a1) at (1*\dx,\ya) {random\\ frames,\\ batch 16};
  \node[tn, n]    (a2) at (2*\dx,\ya) {paired\\ geometry,\\ mosaic on};
  \node[tn, fuse] (a4) at (4*\dx,\ya) {\method\\ image mode,\\ empty state};
  \node[tn, n]    (a5) at (5*\dx,\ya) {detection\\ loss\\ $\mathcal{L}_{\text{det}}$};
  \node[tn, n]    (a6) at (6*\dx,\ya) {MuSGD,\\ AMP, EMA;\\ checkpoint A};
  \draw[arr] (a0) -- (a1); \draw[arr] (a1) -- (a2);
  \draw[arr] (a2) -- node[lbl, above=1pt] {no corruption} (a4);
  \draw[arr] (a4) -- (a5); \draw[arr] (a5) -- (a6);

  % ---------- stage B ----------
  \node[grp, anchor=west] at (-0.8,{\yb+1.2}) {Stage B, clips};
  \node[tn, dat]   (b0) at (0*\dx,\yb) {sequences:\\ RGBT-Tiny,\\ VT-MOT,\\ VisDrone VID};
  \node[tn, n]     (b1) at (1*\dx,\yb) {clips of\\ $L$ = 8 to 16,\\ stride 1--3};
  \node[tn, n]     (b2) at (2*\dx,\yb) {shared\\ geometry\\ per clip};
  \node[tn, alert] (b3) at (3*\dx,\yb) {corruption\\ simulator};
  \node[tn, tim]   (b4) at (4*\dx,\yb) {\methodS\\ unrolled,\\ burn-in 2,\\ TBPTT};
  \node[tn, n]     (b5) at (5*\dx,\yb) {loss\\ $\overline{\mathcal{L}}_{\text{det}}$ $+$\\ $\lambda_{\text{rel}}\,$BCE};
  \node[tn, n]     (b6) at (6*\dx,\yb) {optimiser;\\ FP32 scan;\\ checkpoint B};
  \foreach \i/\j in {0/1,1/2,2/3,3/4,4/5,5/6} \draw[arr] (b\i) -- (b\j);

  % simulator details above its node
  \node[lbl, anchor=south, text width=2.3cm, text=cAlert] at (3*\dx,{\yb+0.82})
       {bursty dropout, blackout,\\ drift, desync, gaps};

  % checkpoint A initialises stage B
  \draw[arr] (a6.south) -- ++(0,-0.5) -| node[lbl, pos=0.22, above=0.5pt] {initialise} ($(b4.north)+(0.35,0)$);

  % reliability targets from the simulator to the loss
  \draw[darr, draw=cAlert] (b3.south) -- ++(0,-0.42) -| (b5.south);
  \node[lbl, anchor=north, text=cAlert] at (4.55*\dx,{\yb-1.25}) {reliability targets $q$};

  % state carried to the next clip
  \draw[darr, draw=cTime, over] ($(b4.south)+(-0.35,0)$) -- ++(0,-1.25) -| (b1.south);
  \node[lbl, anchor=north, text=cTime] at (2.5*\dx,{\yb-2.08}) {final state initialises the next clip, $p=0.5$};

  % ---------- evaluation ----------
  \node[grp, anchor=west] at (-0.8,{\yc+1.15}) {Evaluation};
  \node[base, fuse, fill=cFuse!10, text width=8.4cm, minimum height=1.45cm, font=\scriptsize, align=left] (ev) at (2.45*\dx,\yc)
     {\textbf{image metrics}\enspace AP, AP\textsubscript{50}, AP\textsubscript{S}, AP below 16\,px, SAFit\\[2pt]
      \textbf{stream robustness}\enspace degradation curves, recovery time, flicker\\[2pt]
      \textbf{efficiency}\enspace latency with NMS, memory against sequence length};
  \node[tn, n, minimum height=1.45cm] (log) at (6*\dx,\yc) {run log:\\ seed, config\\ hash, tables};
  \draw[arr] (b6.south) -- ++(0,-1.6) -| ($(ev.north)+(3.4,0)$);
  \draw[arr] (ev) -- (log);
\end{tikzpicture}

\vspace{2mm}
\caption[The two-stage training pipeline]{The two-stage training pipeline.
Stage A trains the still-image \method{} and yields the Phase~1 model.
Stage B continues on clips with a corruption simulator, truncated
back-propagation through time (TBPTT) and stochastic carry-over of the state
between clips. Because the simulator knows which sensor it disabled or
shifted, it provides targets for the reliability heads. Every evaluation is
logged with its seed and configuration hash, and the tables of the final
paper are generated from these logs.}
\label{fig:training}
\end{figure}


\section{Inference modes and complexity}
\label{sec:modes}

The same weights run in three modes.
\begin{itemize}
  \item \emph{Image mode} uses an empty state for each frame. It is used in
        Phase~1 and as an internal baseline.
  \item \emph{Clip mode} processes a stored video causally from its first
        frame. It is used in Phase~2. A bidirectional variant can be run
        offline as an upper bound.
  \item \emph{Stream mode} keeps the state between calls and advances it by
        one step per arriving frame. It is used in Phases~3 and~4.
\end{itemize}
Table~\ref{tab:complexity} compares the per-frame cost of the temporal
options that the ablation will test.

\begin{table}[htbp]
\centering
\caption[Per-frame cost of temporal designs]{Per-frame cost of temporal designs
as a function of the number of frames already seen, $T$, and the window
length, $k$. $M$ is the number of feature locations.}
\label{tab:complexity}
\small
\begin{tabularx}{\textwidth}{@{}P{3.4cm}P{2.3cm}P{2.3cm}L@{}}
\hd{Design} & \hd{Compute} & \hd{Memory} & \hd{Behaviour} \\ \gap
per-frame detector        & $\mathcal{O}(1)$   & none               & no temporal context \\
sliding window of $k$ frames & $\mathcal{O}(k)$ & $\mathcal{O}(kM)$ & context limited to $k$ frames \\
temporal attention over $k$ & $\mathcal{O}(k M)$ & $\mathcal{O}(kM)$ & quadratic in $k$ if all pairs attend \\
convolutional GRU          & $\mathcal{O}(M)$   & $\mathcal{O}(M)$  & unbounded memory; gates not tied to time \\
Temporal State Memory      & $\mathcal{O}(M N)$ & $\mathcal{O}(MN)$ & unbounded memory; step tied to reliability and true interval \\
\end{tabularx}
\end{table}

\section{Hypotheses and how they could fail}
\label{sec:hypotheses}

Each hypothesis is paired with the comparison that tests it and the outcome
that would refute it (Table~\ref{tab:hypotheses}). The targets are
deliberately modest. A refuted hypothesis is a publishable finding if the
experiment is sound, so the protocol of Chapter~\ref{ch:protocol} is designed
to make every outcome interpretable.

\begin{table}[htbp]
\centering
\caption[Hypotheses, tests and refutation criteria]{Hypotheses, the tests that
address them and the outcomes that would refute them. ``Seed spread'' is the
standard deviation over three training seeds.}
\label{tab:hypotheses}
\small
\begin{tabularx}{\textwidth}{@{}P{0.8cm}P{3.6cm}P{3.9cm}L@{}}
\hd{} & \hd{Hypothesis} & \hd{Test} & \hd{Refuted if} \\ \gap
H1 & Temporal memory improves tiny-object detection. & \methodS{} against image-mode \method{} on RGBT-Tiny, AP for objects under 16\,px and SAFit. & The gain is below twice the seed spread. \\ \sgap
H2 & The selective state is the right memory. & TSM against convolutional GRU, a sliding window and windowed temporal attention at matched GFLOPs. & Any alternative matches TSM's accuracy at lower per-frame latency. \\ \sgap
H3 & Gating the step by reliability gives graceful degradation. & Bursty thermal dropout and misalignment, against output gating, COMO-style fusion and concatenation. & Area under the degradation curve is not better, or recovery after an outage is not faster. \\ \sgap
H4 & Scaling the step by the true interval generalises across frame rates. & Train at 15\,FPS without stride augmentation; test at 7.5 and 3.75\,FPS and with random drops, with and without the factor $\tau_t/\tau_0$. & Accuracy loss under rate shift is not reduced. \\ \sgap
H5 & The streaming model is deployable. & TensorRT FP16 export, desktop GPU and Jetson Orin, end-to-end latency including NMS and motion estimation. & Under 30\,FPS on the desktop GPU or under 15\,FPS on Orin at $640\times512$, or more than 1 AP\textsubscript{50} lost in export. \\
\end{tabularx}
\end{table}
